\documentclass[11pt,a4paper]{article}
\usepackage[T1]{fontenc}
\usepackage[utf8]{inputenc}
\usepackage{amsmath,amssymb,xcolor}
\usepackage[margin=25mm]{geometry}
\usepackage[colorlinks=true,urlcolor=blue]{hyperref}
\definecolor{nexorablue}{HTML}{2459B5}
\setlength{\parindent}{0pt}
\setlength{\parskip}{7pt}
\setlength{\emergencystretch}{2em}
\renewcommand{\familydefault}{\sfdefault}
\begin{document}

{\small\bfseries\color{nexorablue} NEXORA CPGE}\par
{\small MP / NEX-IND-06 / Énoncé}\par
{\LARGE\bfseries Après la disparition du flux extérieur}\par
\vspace{0.25em}\hrule\vspace{0.6em}

Lorsque le flux extérieur se fige, une boucle de $R=5{,}0\,\mathrm\Omega$ et $L=0{,}25\,\mathrm H$ porte encore $i_0=0{,}80\,\mathrm A$. Donner le courant et l’énergie magnétique un temps $L/R$ après cet instant. Déterminer l’énergie totale dissipée à partir de cet instant initial.


\end{document}
